\scanpage{110}
donc, pour cette raison, savoir définir les
$\mathbb Rf_*,\ \mathbb R_!f_*$ pour des arguments qui ne sont
pas dans $D^+$, \inserted{comme foncteurs $D^-(X)-D^-(Y),\ D(X)-D(Y),$}
Pour $\mathbb Rf_*$, ce sera possible
dès que $f_*$ est de dim cohomologique finie.
C'est peut-être vrai dès que $Y$ est noethérien,
cela \deleted{est} \inserted{résulte} en tous cas de la finitude de la
dim ch. si $Y$ est, \deleted{soit de} soit de type
fini sur un corps, soit de type fini sur
$\operatorname{Spec}\mathbb Z$. Pour $\mathbb R_!f_*$, on est en meilleure
posture encore, grâce au changement de base propre,
donc \uline{si} $Y$ \uline{qu. cpt.}, on a
\[
\diagram{\mathbb R_!f_*:\quad
\left.\begin{array}{ccc}
D(X)&-&D(Y)\\
D^+(X)&-&D^+(Y)\\
D^-(X)&-&D^-(Y)\\
D^b(X)&-&D^b(Y)
\end{array}\right\}\quad
\begin{gathered}\text{si }Y\text{ q. cpt}\\ f\text{ compactifiable}\end{gathered}}
\]
Comme \unsure{on} y est, notons,
\[
\diagram{\mathbb Rf_*:\quad
\left.\begin{array}{ccc}
D(X)&-&D(Y)\\
D^+(X)&-&D^+(Y)\\
D^-(X)&-&D^-(Y)\\
D^b(X)&-&D^b(Y)
\end{array}\right\}\quad
\begin{gathered}\text{si }Y\text{ de t.f.}\\
\text{sur un corps}\\\text{ou sur }\operatorname{Spec}\mathbb Z,\\
\text{et }f\text{ de t.f.}\end{gathered}}
\]
\[
\diagram{\mathbb Lf^*:\quad
\begin{array}{ccc}
D(\unsure{Y})&-&D(X)\\
D^-(Y)&-&D^-(X)\\
D^+(Y)&-&D^+(X)
\end{array}}
\qquad D^b(Y)\longrightarrow D^b(X)
\]
\mspage{109}
